\cite{OKI-78SR_Series} % Example of a citation to a reference in the bibliography

Reference to Figure~\vref{fig:gallery}. % The \vref command specifies the location of the reference

\begin{figure}[tb]
\centering 
\includegraphics[width=0.5\columnwidth]{GalleriaStampe} 
\caption[An example of a floating figure]{An example of a floating figure (a reproduction from the \emph{Gallery of prints}, M.~Escher,\index{Escher, M.~C.} from \url{http://www.mcescher.com/}).} % The text in the square bracket is the caption for the list of figures while the text in the curly brackets is the figure caption
\label{fig:gallery} 
\end{figure}

\lipsum[13] % Dummy text

\begin{table}[hbt]
\caption{Table of Grades}
\centering
\begin{tabular}{llr}
\toprule
\multicolumn{2}{c}{Name} \\
\cmidrule(r){1-2}
First name & Last Name & Grade \\
\midrule
John & Doe & $7.5$ \\
Richard & Miles & $2$ \\
\bottomrule
\end{tabular}
\label{tab:label}
\end{table}

Reference to Table~\vref{tab:label}. % The \vref command specifies the location of the reference

Reference the figure composed of multiple subfigures as Figure~\vref{fig:esempio}. Reference one of the subfigures as Figure~\vref{fig:ipsum}. % The \vref command specifies the location of the reference

\lipsum[15-18] % Dummy text

\begin{figure}[tb]
\centering
\subfloat[A city market.]{\includegraphics[width=.45\columnwidth]{Lorem}} \quad
\subfloat[Forest landscape.]{\includegraphics[width=.45\columnwidth]{Ipsum}\label{fig:ipsum}} \\
\subfloat[Mountain landscape.]{\includegraphics[width=.45\columnwidth]{Dolor}} \quad
\subfloat[A tile decoration.]{\includegraphics[width=.45\columnwidth]{Sit}}
\caption[A number of pictures.]{A number of pictures with no common theme.} % The text in the square bracket is the caption for the list of figures while the text in the curly brackets is the figure caption
\label{fig:esempio}
\end{figure}
\cite{CMT-8504-100-SMT-TR} % Example of a citation to a reference in the bibliography

\begin{description}
\item[Word] Definition
\item[Concept] Explanation
\item[Idea] Text
\end{description}

\lipsum[12] % Dummy text

\begin{itemize}[noitemsep] % [noitemsep] removes whitespace between the items for a compact look
\item First item in a list
\item Second item in a list
\item Third item in a list
\end{itemize}


\subsection{Introduction}
 
Le projet Vigizone demande une architecture matérielle complexe qui ne peut être réalisée par une ou plusieurs cartes toute faites.
Il est donc nécessaire de concevoir une carte électronique sur mesure pour répondre aux besoins du projet.
Les cartes doivent être capables de gérer les différents capteurs et actionneurs nécessaires au fonctionnement du système, tout en respectant les contraintes de taille, de consommation d'énergie et de coût.
Pour se faire, nous nous sommes appuyés sur le cahier des charges réalisé à l'aide de GTM OUEST, pour la conception technique.

%------------------------------------------------

\subsection{Choix des composants}

\lipsum[11] % Dummy text

\subsection{Hub}

\lipsum[12] % Dummy text


\subsection{Ancre}


%------------------------------------------------

\subsection{Ouvrier}


\subsection{réalisation PCB}

\subsection{Soudure PCB}

\subsection{Test Physique}

\subsection{Résultat final}